\chapter{Dictionary of sign conventions}
\label{app:sign-conventions}


Sign conventions used in different sources are compared and translated below.

\section{Loday--Vallette vs. this manuscript}

\begin{center}
\begin{longtable}{|>{\raggedright\arraybackslash}p{5.8cm}|>{\raggedright\arraybackslash}p{3.9cm}|>{\raggedright\arraybackslash}p{3.9cm}|}
\caption{Sign convention comparison} \\
\hline
\textbf{Object/Operation} & \textbf{Loday--Vallette \cite{LV12}} & \textbf{This Manuscript} \\
\hline
\endfirsthead
\multicolumn{3}{c}{\tablename\ \thetable\ -- \textit{Continued}} \\
\hline
\textbf{Object/Operation} & \textbf{Loday--Vallette} & \textbf{This Manuscript} \\
\hline
\endhead
\hline
\endfoot
\hline
\endlastfoot
Koszul sign rule & $(-1)^{|a| \cdot |b|}$ & $(-1)^{|a| \cdot |b|}$ (same) \\
Differential degree & $|d| = -1$ & $|d| = +1$ \\
Suspension & $s: V \to sV$, $|sv| = |v| + 1$ & $s: V \to V[-1]$, $|sv| = |v| + 1$ (cohomological: $V[-1]^k = V^{k-1}$, so $v \in V^m$ maps to $sV^{m+1}$) \\
Desuspension & $s^{-1}: sV \to V$, $|s^{-1}(sv)| = |v|$ & $[-1]: V[1] \to V$ \\
Bar construction & $B(\mathcal{A}) = (T^c(s\mathcal{A}), d_{\text{bar}})$ & $\bar{B}(\mathcal{A}) = (T^c(s^{-1}\bar{\mathcal{A}}), d)$ (desuspension; cf.\ cohomological convention) \\
Cobar construction & $\Omega(\mathcal{C}) = (T(s^{-1}\mathcal{C}), d_{\text{cobar}})$ & $\Omega(\mathcal{C})_n = \mathcal{C}^{\otimes_{\text{co}} n}$ \\
Coproduct sign & $\Delta(ab) = \Delta(a) \cdot \Delta(b)$ with $(-1)^{|a_2| \cdot |b_1|}$ & Same \\
Shuffle sign & $(a \shuffle b) = \sum (-1)^{\sigma} a_{\sigma(1)} \otimes \cdots$ & Same \\
Collision divisor ordering & Not applicable (no geometry) & Lexicographic: $i < j$ always \\
Logarithmic form sign & Not applicable & $\eta_{ji} = \eta_{ij}$ as differential forms ($d\log(-1)=0$). Note: the Orlik--Solomon generators $\omega_{ij}$ are antisymmetric ($\omega_{ji} = -\omega_{ij}$); the identification $\omega_{ij} \mapsto \eta_{ij}$ for $i < j$ fixes an orientation. \\
Arnold relations & Not applicable & $\sum_{\text{cyclic}} \eta_{ij} \wedge \eta_{jk} = 0$ \\
\end{longtable}
\end{center}

\subsection{Key differences explained}

\emph{1.~Differential degree.}
Loday--Vallette use homological grading ($|d| = -1$); this manuscript uses cohomological grading ($|d| = +1$).

\emph{Translation.} If $(C_*, d)$ is a chain complex in LV convention ($d: C_n \to C_{n-1}$), set $C^n := C_{-n}$ with the same underlying map~$d$, which now raises cohomological degree by~$1$.  No sign change is needed: $d_{\text{us}} = d_{\text{LV}}$ (as linear maps).

\emph{2.~Suspension.}
Loday--Vallette treat suspension $s$ as an explicit operator; this manuscript uses the degree shift $sV = V[-1]$, so $|sv| = |v| + 1$.

\emph{Translation.} $sV$ (LV) corresponds to $V[-1]$ (us), i.e., our $sV$.  Both raise degree by~$1$: LV's convention $(sV)_k = V_{k-1}$ translates via $\deg_{\mathrm{coh}} = -\deg_{\mathrm{hom}}$ to $(sV)^k = V^{k-1} = (V[-1])^k$.

\emph{3.~Bar/cobar formulas.}
Loday--Vallette include explicit suspension in the formula; this manuscript absorbs it into the D-module structure.

\emph{Translation.} Our $\bar{B}_n(\mathcal{A})$ corresponds to LV's $B(\mathcal{A})$ in degree $n$ after removing the suspension $s$.

\section{Beilinson--Drinfeld vs. this manuscript}

\begin{center}
\begin{tabular}{|>{\raggedright\arraybackslash}p{5.8cm}|>{\raggedright\arraybackslash}p{3.9cm}|>{\raggedright\arraybackslash}p{3.9cm}|}
\hline
\textbf{Object/Operation} & \textbf{BD \cite{BD04}} & \textbf{This Manuscript} \\
\hline
Chiral algebra & $\mathcal{A}$ (D-module) & $\mathcal{A}$ (same) \\
Factorization & $j_{!*}\mathcal{A}$ & Implicit in $\ConfigSpace{n}$ \\
Configuration space & $\text{Ran}(X)$ & $\bigcup_n \ConfigSpace{n}$ \\
Collision divisor & $(i,j)^c$ (complement) & $\omitpair{i}{j}$ (hat notation) \\
Chiral connection & $\nabla$ & $\nabla$ (same) \\
Residue & Implicit in factorization & Explicit via $\text{Res}_{z_i \to z_j}$ \\
OPE & Encoded in $\mathcal{A}(D)$ & Explicit: $a(z)b(w) = \sum c_n(w)/(z-w)^n$ \\
Koszul duality & Not discussed (BD focus on one side) & Central theme \\
\hline
\end{tabular}
\end{center}

\subsection{Key differences explained}

\emph{BD's approach.}
\begin{itemize}
\item Ran space formulation (taking colimits over all $n$)
\item Factorization axiom as the primary structure
\item D-modules as the fundamental objects
\item No explicit bar/cobar constructions
\end{itemize}

\emph{This monograph.}
\begin{itemize}
\item Explicit configuration spaces $\ConfigSpace{n}$ for each $n$
\item Bar complex as an explicit resolution
\item Koszul duality as the organizing principle
\item Geometric realization via residues
\end{itemize}

\emph{Translation.} Our bar complex $\bar{B}_*(\mathcal{A})$ is the ``Chevalley--Cousin resolution'' of $\mathcal{A}$ viewed as a factorization algebra in BD's sense (BD Theorem 3.4.9).

\section{Costello--Gwilliam vs. this manuscript}

\begin{center}
\begin{tabular}{|>{\raggedright\arraybackslash}p{5.8cm}|>{\raggedright\arraybackslash}p{3.9cm}|>{\raggedright\arraybackslash}p{3.9cm}|}
\hline
\textbf{Object/Operation} & \textbf{CG \cite{CG17}} & \textbf{This Manuscript} \\
\hline
Factorization algebra & $\mathcal{F}$ (precosheaf) & $\mathcal{A}$ (chiral algebra) \\
Configuration space & $\text{Conf}_n(M)$ & $\ConfigSpace{n}$ \\
Compactification & Fulton--MacPherson & Same \\
Bar complex & Not emphasized & Central construction \\
Feynman diagrams & Graph complex & Genus-stratified bar complex \\
Renormalization & BV formalism & Completion of bar complex \\
Quantum corrections & Loop order & Genus $g$ \\
\hline
\end{tabular}
\end{center}

\subsection{Key differences explained}

\emph{CG's approach.}
\begin{itemize}
\item Factorization algebras on manifolds (any dimension)
\item Feynman diagram expansion for observables
\item BV formalism for quantization
\item Emphasis on renormalization and effective field theory
\end{itemize}

\emph{This monograph.}
\begin{itemize}
\item Chiral algebras on curves (1-dimensional)
\item Genus expansion instead of loop expansion
\item Bar-cobar construction and Koszul duality instead of BV
\item Emphasis on algebraic structure and representation theory
\end{itemize}

\emph{Translation.}
\begin{itemize}
\item CG's ``factorization algebra'' on a curve $X$ = BD/Our ``chiral algebra''
\item CG's ``observable'' = Our ``element of $\mathcal{A}$''
\item CG's ``1-loop graph'' = Our ``genus-1 contribution to bar complex''
\item CG's ``renormalization'' = Our ``nilpotent completion of bar complex''
\end{itemize}

\section{Kontsevich vs. this manuscript}

\begin{center}
\begin{tabular}{|>{\raggedright\arraybackslash}p{5.8cm}|>{\raggedright\arraybackslash}p{3.9cm}|>{\raggedright\arraybackslash}p{3.9cm}|}
\hline
\textbf{Object/Operation} & \textbf{Kontsevich \cite{Kon99}} & \textbf{This Manuscript} \\
\hline
Configuration space & $C_n(\mathbb{R}^d)$ (open) & $\ConfigSpace{n}$ (compactified) \\
Forms & Smooth forms & Logarithmic forms \\
Propagator & $\frac{1}{z-w}$ & $\frac{dz}{z-w}$ (with log pole) \\
Graphs & Directed graphs & Genus-stratified graphs \\
Formality & $\mathcal{U}: T_{\text{poly}} \to D_{\text{poly}}$ & Bar-cobar: $\bar{B} \dashv \Omega$ \\
Wheels & Non-planar graphs & Higher genus contributions \\
\hline
\end{tabular}
\end{center}

\subsection{Key differences explained}

\emph{Kontsevich's formality.}
\begin{itemize}
\item Classical: Poisson manifolds $\to$ deformation quantization
\item Configuration spaces in $\mathbb{R}^d$ (no compactification)
\item Smooth forms (no poles)
\item Graph complex (combinatorial)
\end{itemize}

\emph{Our chiral formality.}
\begin{itemize}
\item Coisson algebras (PVA) $\to$ vertex algebras (OPE) via chiral configuration spaces
\item Configuration spaces on curves (compactified)
\item Logarithmic forms (poles at collisions)
\item Genus-stratified complex (geometric)
\end{itemize}

\emph{Relation.} Kontsevich's formality is the \emph{genus-0 part} of chiral formality. The bar-cobar framework extends Kontsevich to include all genera (higher-loop corrections).

\section{Summary table: all conventions}

\begin{center}
\begin{longtable}{|>{\raggedright\arraybackslash}p{5cm}|>{\raggedright\arraybackslash}p{2cm}|>{\raggedright\arraybackslash}p{2cm}|>{\raggedright\arraybackslash}p{2cm}|>{\raggedright\arraybackslash}p{2cm}|}
\caption{Sign convention dictionary} \\
\hline
\textbf{Operation} & \textbf{LV} & \textbf{BD} & \textbf{CG} & \textbf{Us} \\
\hline
\endfirsthead
\multicolumn{5}{c}{\tablename\ \thetable\ -- \textit{Continued}} \\
\hline
\textbf{Operation} & \textbf{LV} & \textbf{BD} & \textbf{CG} & \textbf{Us} \\
\hline
\endhead
Koszul sign & $(-1)^{|a||b|}$ & $(-1)^{|a||b|}$ & $(-1)^{|a||b|}$ & $(-1)^{|a||b|}$ \\
Differential degree & $-1$ & $+1$ & $+1$ & $+1$ \\
Suspension degree & $+1$ & $+1$ & $+1$ & $+1$ \\
Collision sign & N/A & Implicit & Explicit & Explicit \\
Arnold relations & N/A & Implicit & Yes & Yes \\
\hline
\end{longtable}
\end{center}

\emph{Recommendation.} When reading other sources, always consult this appendix to translate signs and conventions. The most common source of error in chiral algebra computations is inconsistent sign conventions.

\section{Practical guide: how to translate}

\subsection{From Loday--Vallette to our conventions}

\emph{Rule~1.} Replace $|d| = -1$ with $|d| = +1$.

\emph{Rule~2.} Replace LV's suspension $s$ with our degree shift $[-1]$: $sV = V[-1]$.

\emph{Rule~3.} Our geometric forms $\Omega^n(\log D)$ correspond to LV's suspended generators $s^n V$.

\subsection{From Beilinson--Drinfeld to our conventions}

\emph{Rule~1.} BD's $j_{!*}$ (factorization pushforward) = our explicit configuration space integral.

\emph{Rule~2.} BD's Ran space $\text{Ran}(X)$ = our $\bigcup_n \ConfigSpace{n} / \sim$ (colimit).

\emph{Rule~3.} BD's ``chiral product'' = our ``OPE residue''

\subsection{From Costello--Gwilliam to our conventions}

\emph{Rule~1.} CG's ``1-loop'' = our ``genus 1''

\emph{Rule~2.} CG's ``renormalization'' = our ``$I$-adic completion''

\emph{Rule~3.} CG's BV bracket = our bar differential (when specialized to curves).

\section{Examples of translation}

\begin{example}[Translating a bar differential formula from LV]
\emph{Loday--Vallette writes.}
\[d_{LV}(s a_1 \otimes s a_2 \otimes s a_3) = (-1)^{|a_1|} (s[a_1, a_2] \otimes s a_3) + \cdots\]

\emph{Our notation.}
\[d(a_1 \otimes a_2 \otimes a_3 \otimes \omega_2) = ([a_1, a_2] \otimes a_3 \otimes \text{Res}[\omega_2]) + \cdots\]

\emph{Translation.}
\begin{itemize}
\item Remove explicit suspension $s$ (implicit in form degree)
\item Change sign convention: $(-1)^{|a_1|}$ becomes automatic from Koszul rule
\item Add geometric form $\omega_2 \in \Omega^2(\log D)$
\end{itemize}
\end{example}

\begin{example}[Translating a BD factorization formula]
\emph{Beilinson--Drinfeld writes.}
\[j_{!*}(\mathcal{A} \boxtimes \mathcal{A})|_{U \sqcup V} \simeq j_{!*}\mathcal{A}|_U \boxtimes j_{!*}\mathcal{A}|_V\]

\emph{Our notation.}
\[\barcomplex{1}{\mathcal{A}}|_{U \sqcup V} = \int_{C_2(U) \sqcup C_2(V)} \mathcal{A}^{\boxtimes 2} \otimes \Omega^1_{\log} \simeq \barcomplex{1}{\mathcal{A}}|_U \otimes \barcomplex{1}{\mathcal{A}}|_V\]

\emph{Translation.}
\begin{itemize}
\item $j_{!*}$ (pushforward) = configuration space integral
\item Factorization isomorphism = bar complex splitting
\item Add explicit forms $\Omega^1_{\log}$
\end{itemize}
\end{example}

\section{Sign rules for this manuscript}

For easy reference, here are ALL sign rules used in this manuscript in one place:

\subsection{Koszul signs}

\emph{Rule.} When permuting two graded objects $a$ and $b$:
\[a \otimes b = (-1)^{|a| \cdot |b|} b \otimes a\]

\emph{Application.} Moving a form $\omega$ of degree $|\omega| = n$ past fields $\phi_1, \ldots, \phi_k$:
\[(\phi_1 \otimes \cdots \otimes \phi_k) \otimes \omega = (-1)^{n(|\phi_1| + \cdots + |\phi_k|)} \omega \otimes (\phi_1 \otimes \cdots \otimes \phi_k)\]

\subsection{Collision divisor signs}

\emph{Rule.} Collision divisors are always written with indices in increasing order: $\colldiv{i}{j}$ means $i < j$.

\emph{Convention.} We write $\colldiv{i}{j}$ with $i < j$.  Geometrically $D_{ij} = D_{ji}$ (same subvariety $z_i = z_j$).  The sign convention $\colldiv{j}{i} = -\colldiv{i}{j}$ applies to the \emph{oriented divisor} (i.e., to the associated residue forms, not the subvariety itself).

\subsection{Arnold relation signs}

\emph{3-term Arnold relation.}
\[\LogForm{i}{j} \wedge \LogForm{j}{k} + \LogForm{j}{k} \wedge \LogForm{k}{i} + \LogForm{k}{i} \wedge \LogForm{i}{j} = 0\]

\emph{Higher-degree consequences.}
All higher-degree relations in the OS algebra follow from the 3-term Arnold relations by repeated application. There is no independent ``general Arnold relation'' beyond the 3-term version.

\subsection{Residue signs}

The sign arising from the residue at $\colldiv{i}{j}$ is determined by the Koszul sign rule: the residue operation (cohomological degree~$1$) must commute past the preceding tensor factors.  For elements $s^{-1}a_1 \otimes \cdots \otimes s^{-1}a_n$ in the bar complex, the residue at $D_{ij}$ (colliding factors $i$ and $j$) picks up the sign $(-1)^{\sum_{k<i}|s^{-1}a_k|}$.  In particular, when all generators have even desuspended degree, all residue signs are positive.

For the explicit sign computations relevant to the bar differential, see the proof of $d^2 = 0$ in \S\ref{sec:bar-nilpotency}.

\section{Common pitfalls and how to avoid them}

\subsection{Pitfall 1: forgetting Koszul signs}

\emph{Wrong.} $d(a \otimes b \otimes \omega) = da \otimes b \otimes \omega + a \otimes db \otimes \omega$

\emph{Right.} $d(a \otimes b \otimes \omega) = da \otimes b \otimes \omega + (-1)^{|a|} a \otimes db \otimes \omega$

The sign $(-1)^{|a|}$ comes from moving $d$ (degree +1) past $a$ (degree $|a|$).

\subsection{Pitfall 2: confusing hat notations}

\emph{Wrong.} $\widehat{ij}$ could mean ``omit factors $i$ and $j$,'' or ``the set $\{i, j\}$,'' or something else entirely.

\emph{Right.} Use our convention: $\omitpair{i}{j}$ always means ``omit both $i$ and $j$'' with no ambiguity.

\subsection{Pitfall 3: collision divisor ordering}

\emph{Convention.} We write collision divisors with $i < j$: $\colldiv{1}{3}$ rather than $\colldiv{3}{1}$.  Since $\colldiv{i}{j} = \{z_i = z_j\}$ is an unordered locus, $\colldiv{3}{1} = \colldiv{1}{3}$ (no sign).  Likewise, the logarithmic forms satisfy $\LogForm{j}{i} = \LogForm{i}{j}$ since $d\log(z_j - z_i) = d\log(-(z_i - z_j)) = d\log(z_i - z_j)$.

The sign convention $e_{ji} = -e_{ij}$ applies only to Orlik--Solomon algebra generators (an algebraic convention for exterior algebra presentations), not to the geometric forms or divisors.

\subsection{Pitfall 4: Arnold relation orientation}

\emph{Wrong.} Assuming Arnold relations hold without signs.

\emph{Right.} The cyclic sum includes signs from wedge product antisymmetry:
\[\LogForm{1}{2} \wedge \LogForm{2}{3} + \LogForm{2}{3} \wedge \LogForm{3}{1} + \LogForm{3}{1} \wedge \LogForm{1}{2} = 0\]

Note: Each term has a $+$ sign in the cyclic sum, but the individual forms anticommute.

\section{Summary and recommendations}

\emph{For readers unfamiliar with sign conventions.}
\begin{enumerate}
\item Start with Remark~\ref{rem:LV-signs} (Loday--Vallette comparison, in Chapter~\ref{chap:bar-cobar})
\item Consult this appendix when reading other sources
\item Work through the examples in \S\ref{sec:heisenberg-bar-complex} (Heisenberg explicit calculations)
\item Verify signs match in low-degree computations
\end{enumerate}

\emph{For experts.}
Our conventions match:
\begin{itemize}
\item Kontsevich's graph complex (lexicographic ordering)
\item Costello--Gwilliam's factorization algebras (cohomological grading)
\item Beilinson--Drinfeld's D-module perspective (implicit signs)
\end{itemize}

The main difference from Loday--Vallette is grading direction ($|d| = +1$ vs $-1$); the translation is a global sign reversal on the grading.

\section{Sign convention dictionary}
\label{sec:sign-dictionary-complete}

\begin{table}[ht!]
\centering
\caption{Sign conventions across sources}
\label{tab:sign-conventions-complete}
\small
\begin{tabular}{|l|c|c|c|c|}
\hline
\textbf{Item} & \textbf{Ours} & \textbf{BD} & \textbf{CG} & \textbf{LV} \\
\hline
\hline
Differential degree & $|d| = +1$ & implicit & $+1$ & $-1$ \\
Koszul sign rule & $(-1)^{|a||b|}$ & yes & yes & yes \\
Suspension & $s: V \to V[-1]$ & yes & yes & opposite \\
Bar complex degree & increasing & yes & yes & decreasing \\
Logarithmic 1-forms & $\eta_{ij} = \frac{dz_i - dz_j}{z_i - z_j}$ & yes & yes & N/A \\
Residue sign & $\text{Res}_{z=w}[\eta_{ij}] = +1$ & yes & yes & N/A \\
Factorization & left-to-right & yes & opposite & N/A \\
Collision divisor & $D_{ij}$ with $i < j$ & yes & same & N/A \\
Transposition & $\eta_{ji} = \eta_{ij}$ (forms) & yes & yes & N/A \\
Wedge product & $a \wedge b = (-1)^{|a||b|} b \wedge a$ & yes & yes & yes \\
Verdier duality & $\mathbb{D}(\mathcal{F}) = R\mathcal{H}om(\mathcal{F}, \omega[d])$ & yes & same & N/A \\
Central extension & $[a,b] = ab - ba$ & yes & yes & yes \\
A-infinity & $\mu_n: \mathcal{A}^{\otimes n} \to \mathcal{A}$ & implicit & yes & yes \\
\hline
\end{tabular}
\end{table}

\emph{Key to abbreviations.}
\begin{itemize}
\item BD = Beilinson--Drinfeld \cite{BD04}
\item CG = Costello--Gwilliam \cite{CG17}
\item LV = Loday--Vallette \cite{LV12}
\end{itemize}

\subsection{Conversion formulas}

\begin{proposition}[Loday--Vallette conversion; \ClaimStatusProvedHere]\label{prop:LV-conversion-complete}
To convert from our conventions to Loday--Vallette:
\begin{enumerate}
\item Replace $|d| = +1$ with $|d| = -1$ (flip grading direction)
\item Replace bar complex $\bar{B}^n$ with $\bar{B}_{-n}$ (decreasing filtration)
\item Keep Koszul sign rule unchanged (all four conventions agree on the Koszul sign rule)
\item Keep composition order unchanged: $(f \circ g)_{\text{LV}} = f \circ g$ (the degree reversal affects signs in the Leibniz rule, not the order of function composition)
\end{enumerate}
\end{proposition}

\begin{proof}[Verification]
The key identity is that homological grading (LV) relates to cohomological grading 
(ours) by $\deg_{\text{hom}} = -\deg_{\text{coh}}$.

For a differential $d$:
\begin{itemize}
\item Our convention: $d: C^n \to C^{n+1}$, so $|d| = +1$
\item LV convention: $d: C_n \to C_{n-1}$, so $|d|_{\text{LV}} = -1$
\end{itemize}

The Koszul sign rule in the Leibniz formula is the same in both conventions:
\[d(a \otimes b) = da \otimes b + (-1)^{|a|} a \otimes db\]
(The degree $|a|$ is taken in whichever grading is active.)  The sign difference between conventions manifests in the \emph{bar differential}, where LV uses explicitly suspended elements $sa$: the bar differential sign $(-1)^{|sa|} = (-1)^{|a|+1}$ in LV's homological convention (where $|sa| = |a|+1$) differs from $(-1)^{|a|}$ in ours.
\end{proof}

\subsection{Explicit sign calculations}

\begin{example}[Three-point Arnold relation]\label{ex:arnold-sign-explicit-dict}
For three points $z_1, z_2, z_3$, the Arnold relation is:
\[\eta_{12} \wedge \eta_{23} + \eta_{23} \wedge \eta_{31} + \eta_{31} \wedge \eta_{12} = 0\]

\emph{Sign verification.}
\begin{enumerate}
\item Each $\eta_{ij}$ has degree 1
\item Wedge product: $\eta \wedge \eta'$ has degree 2
\item Moving forms: $\eta \wedge \eta' = (-1)^{1 \cdot 1} \eta' \wedge \eta = -\eta' \wedge \eta$
\end{enumerate}

All signs are explicit and match BD conventions.
\end{example}

\begin{remark}[Fermionic vs bosonic]\label{rem:fermionic-bosonic-signs-dict}
For fermionic fields $\psi$: $\psi(z)\psi(w) = -\psi(w)\psi(z)$ (Koszul sign with $|\psi| = 1$).

For bosonic fields $\phi$: $\phi(z)\phi(w) = \phi(w)\phi(z)$ (degree $|\phi| = 0$).

The bar complex treats both uniformly via suspension $s: V \to V[-1]$.
\end{remark}

% ================================================================
% PATCH 015: COMPREHENSIVE SIGN DICTIONARY EXPANSION
% ================================================================

\section{Master comparison table}
\label{sec:master-comparison-signs}

\begin{table}[H]
\centering
\caption{Sign conventions across four sources}
\renewcommand{\arraystretch}{1.6}
\footnotesize
\begin{tabular}{|l|c|c|c|c|}
\hline
\textbf{Convention} & \textbf{BD} & \textbf{CG} & \textbf{LV} & \textbf{Us} \\
\hline
\hline
\textbf{Koszul sign} & $(-1)^{|a||b|}$ & $(-1)^{|a||b|}$ & $(-1)^{|a||b|}$ & $(-1)^{|a||b|}$ \\
\hline
\textbf{Bar differential} & $d + d'$ & $d_{\text{int}} + d_{\text{ext}}$ & $d_1$ & $d_{\text{int}} + d_{\text{res}}$ \\
\hline
\textbf{Cobar degree} & $[n]$ & $[n-2]$ & $[-n]$ & $[2-n]$ \\
\hline
\textbf{Suspension} & $s: V \to V[-1]$ & $s: V \to V[-1]$ & $s: V \to V[1]$ & $s: V \to V[-1]$ \\
\hline
\textbf{Desuspension} & $s^{-1}: V \to V[1]$ & $s^{-1}: V \to V[1]$ & $s^{-1}: V \to V[-1]$ & $s^{-1}: V \to V[1]$ \\
\hline
\textbf{Twist map} & $(-1)^{|a||b|}(b \otimes a)$ & $(-1)^{|a||b|}(b \otimes a)$ & $(-1)^{|a||b|}(b \otimes a)$ & $(-1)^{|a||b|}(b \otimes a)$ \\
\hline
\textbf{Differential degree} & $+1$ & $+1$ & $-1$ & $+1$ \\
\hline
\textbf{Coalgebra comult} & $\Delta$ & $\Delta$ & $\Delta^!$ & $\Delta$ \\
\hline
\textbf{Residue orientation} & counterclockwise & contour dep. & N/A & counterclockwise \\
\hline
\textbf{$\hbar$ convention} & not used & $\hbar = 1$ & not used & $\hbar$ explicit \\
\hline
\end{tabular}
\end{table}

\subsection{Detailed conversion formulas}

\subsubsection{Koszul signs}

\begin{definition}[Koszul sign rule]
When permuting graded elements $a$ and $b$ in a tensor product:

\emph{BD, CG, Us convention.}
\[a \otimes b \leftrightarrow b \otimes a \quad \text{with sign} \quad (-1)^{|a| \cdot |b|}\]

\emph{LV convention.}
\[a \otimes b \leftrightarrow b \otimes a \quad \text{with sign} \quad (-1)^{|a| \cdot |b|} \quad \text{(same Koszul rule)}\]

\emph{Key difference.}
LV works in homological grading ($|d| = -1$), using suspended elements $sa$ with $|sa| = |a| + 1$. The apparent sign differences in bar complex formulas arise from the shifted degrees of $sa$, not from a different Koszul rule. Specifically, swapping $sa \otimes sb$ gives $(-1)^{|sa||sb|} = (-1)^{(|a|+1)(|b|+1)}$, which differs from $(-1)^{|a||b|}$ by the factor $(-1)^{|a|+|b|+1}$.
\end{definition}

\begin{example}[Koszul sign in bar construction]
For $a \in \mathcal{A}^m$, $b \in \mathcal{A}^n$, $c \in \mathcal{A}^p$:

\emph{Us/BD/CG.}
\[d(a \otimes b \otimes c) = d(a) \otimes b \otimes c + (-1)^m a \otimes d(b) \otimes c 
+ (-1)^{m+n} a \otimes b \otimes d(c)\]

\emph{LV (on suspended elements $sa, sb, sc$).}
\[d(sa \otimes sb \otimes sc) = d(sa) \otimes sb \otimes sc + (-1)^{m+1} sa \otimes d(sb) \otimes sc
+ (-1)^{m+n+2} sa \otimes sb \otimes d(sc)\]

\emph{Explanation.} The sign differences arise because $|sa| = m+1$, $|sb| = n+1$, $|sc| = p+1$, and the standard Koszul rule is applied to these shifted degrees. The LV Koszul sign rule itself is the same $(-1)^{|x||y|}$.
\end{example}

\subsection{Quick translation table}

\begin{table}[H]
\centering
\caption{Translation formulas between conventions}
\begin{tabular}{|l|c|}
\hline
\textbf{To convert from} & \textbf{Apply this transformation} \\
\hline
\hline
LV $\to$ Us/BD/CG & Multiply internal terms by $(-1)^{|a|}$, add chiral product \\
\hline
CG $\to$ Us & Shift notation: CG's $V[1]$ (convention $V[n]^k = V^{k-n}$) equals our $V[-1]$ \\
\hline
BD $\to$ Us & Split $d''$ into $d_{\text{residue}}$ and $d_{\text{correction}}$ \\
\hline
Us $\to$ BD & Combine $d_{\text{residue}} + d_{\text{correction}}$ into $d''$ \\
\hline
Us $\to$ LV & Remove chiral product, multiply by $(-1)^{|a|}$ \\
\hline
\end{tabular}
\end{table}

\subsection{Recommendations for readers}

\begin{itemize}
\item \emph{Reading BD.} Our conventions match closely; main difference is we separate out
quantum corrections explicitly.
\item \emph{Reading CG.} The conventions match exactly; use their formulas directly.
\item \emph{Reading LV.} Add extra $(-1)^{|a|}$ signs and remember to include chiral product
which LV does not have.
\item \emph{Writing papers.} We recommend the CG conventions (used here) as they are most explicit about quantum corrections.
\end{itemize}

\begin{remark}[Origins of different conventions]
The conventions reflect different origins: BD developed for D-modules (orientation from holomorphic structure), CG optimized for BV formalism and factorization algebras, LV for pure operadic theory (no geometric residues), and ours emphasizing geometric realization and explicit calculations.
All conventions are \emph{mathematically equivalent}---they are different presentations of the same underlying structures.
\end{remark}
